Agenda Pro — Documento de Escopo MVP v1.0
1. Visão Geral do Produto
O Agenda Pro é um aplicativo mobile para profissionais autônomas (manicures, lash designers, confeiteiras, entre outras) que precisam organizar seus agendamentos e manter comunicação com clientes pelo WhatsApp — sem mudar a forma como trabalham.
O produto resolve três dores centrais:
Perder tempo respondendo mensagens e anotando na agenda de papel
Esquecer de enviar lembretes para as clientes antes dos atendimentos
Não ter aviso antecipado sobre os próprios compromissos do dia
Informações do Produto
Campo
Informação
Produto
Agenda Pro — Aplicativo Mobile
Plataforma
iOS e Android (React Native ou Flutter)
Versão
MVP v1.0
Público-alvo
Profissionais autônomas de beleza e serviços pessoais
Modelo
Freemium — Gratuito com limite + Plano Pro R$29/mês


2. Problema que Resolve
A Ana, manicure autônoma com 8 atendimentos por dia, hoje gerencia tudo manualmente: responde no WhatsApp, anota na agenda de papel e ainda precisa lembrar de mandar mensagem na véspera. São 8 conversas separadas, 8 anotações e muita coisa caindo no esquecimento.
O que o Agenda Pro muda na rotina dela
Cria agendamentos em menos de 30 segundos
Recebe notificação no celular antes de cada atendimento, no intervalo que ela mesma configurou
Envia lembrete para a cliente com um toque, sem digitar nada
Mantém o controle — o envio é feito por ela, mas sem esforço

3. Funcionalidades do MVP
3.1 Cadastro e Onboarding
Cadastro com nome, tipo de serviço e serviços oferecidos com preços
Onboarding em menos de 5 minutos
Configuração do intervalo de notificações pessoais (1 dia antes, 2 horas antes, ou ambos)

3.2 Agenda Principal
Tela principal com vista de dia e semana
Lista de horários em ordem cronológica, estilo agenda de papel
Criação de agendamento: selecionar dia, hora, serviço e nome da cliente
Edição e cancelamento de agendamentos

3.3 Notificações para a Profissional
Funcionalidade central do MVP. A profissional configura, nas preferências, com quanto tempo de antecedência quer ser avisada sobre cada atendimento.
Push notification local — sem servidor, funciona offline
Configuração global: 1 dia antes, 2 horas antes, ou ambos
Ao tocar na notificação, o app abre direto no agendamento correspondente
A tela do agendamento já exibe o botão de envio de lembrete para a cliente — um fluxo único
Exemplo de notificação
“Amanhã às 10h — Manicure com Maria Clara”
Ao tocar: abre o agendamento + botão para enviar lembrete para a Maria Clara.

3.4 Lembretes para a Cliente via WhatsApp
O MVP utiliza dois caminhos de envio, definidos pelo plano da profissional:
Plano Gratuito — Link Manual
O app gera a mensagem pronta e abre o WhatsApp da profissional com o texto já preenchido. Ela só toca em enviar.
Sem API, sem custo por mensagem
Funciona com o número pessoal da profissional
A profissional mantém o controle — zero risco de envio errado
Tecnologia: link wa.me nativo do WhatsApp

Plano Pro — WhatsApp Business API
A profissional conecta um número Business ao app. As mensagens são enviadas automaticamente no horário certo, sem nenhuma ação necessária.
Aprovação de templates junto à Meta
150 mensagens/mês incluídas no plano
Envio automático — zero ação necessária da profissional
Número dedicado, separado do pessoal

Templates incluídos no MVP
Confirmação de agendamento
“Oi [nome]! Tudo certo, seu horário de [serviço] está confirmado para [dia] às [hora]. Qualquer coisa é só chamar!”
Lembrete 24h antes
“Oi [nome]! Passando pra lembrar do seu [serviço] amanhã às [hora]. Te espero!”
Confirmação de presença
“Perfeito! Te aguardo amanhã.”
No plano Pro, todos os templates são editáveis. O app salva a versão personalizada como padrão.

4. Estrutura de Planos
Funcionalidade
Gratuito
Pro — R$29/mês
Agendamentos por mês
Até 30
Ilimitados
Clientes cadastrados
Até 20
Ilimitados
Notificações para a profissional
Configurável
Configurável
Lembrete via link WhatsApp (manual)
Incluído
Incluído
Mensagens automáticas (WhatsApp API)
—
150 msg/mês incluídas
Templates editáveis
2 fixos
Ilimitados
Histórico de clientes
—
Incluído
Suporte prioritário
—
Incluído

Lógica de Upgrade
A profissional começa grátis, usa no dia a dia e quando atingir o limite ou quiser automação total, o upgrade para o Pro acontece de forma natural — ela já percebeu o valor do produto.

5. Telas do MVP
Tela
Descrição
Onboarding
Cadastro de nome, tipo de serviço, serviços e preços. Configuração inicial das notificações pessoais
Agenda — Vista de Dia
Lista de horários do dia em ordem. Botão de adicionar agendamento
Agenda — Vista de Semana
Visão semanal compacta para identificar dias cheios e lacunas
Criar Agendamento
Selecionar data, hora, serviço e nome da cliente
Detalhe do Agendamento
Ver dados do horário. Botões: Enviar lembrete / Editar / Cancelar
Perfil e Configurações
Editar serviços, preços e intervalo de notificações pessoais
Upgrade para Pro
Tela de comparação de planos com call-to-action para assinatura


6. Fora do Escopo do MVP
Itens intencionalmente excluídos para garantir velocidade de lançamento e validação:
Relatórios de receita por serviço ou horário
Sugestões automáticas de reengajamento
Agendamento online pela cliente (self-booking)
Múltiplos profissionais no mesmo app
Integração com pagamentos (Pix, cartão)
Templates de cobrança e pós-atendimento

7. Métricas de Sucesso do MVP
Métrica
Objetivo
Ativação
Profissional cria ao menos 3 agendamentos na primeira semana
Retenção D7
Mais de 50% das usuárias abrem o app na segunda semana
Conversão
Ao menos 5% fazem upgrade para Pro em 30 dias
Satisfação
NPS acima de 40 após 30 dias de uso
Hábito
Usuária ativa abre o app pelo menos 4 vezes por semana


8. Principais Riscos e Mitigações
Risco
Mitigação
A profissional não cria o hábito de abrir o app
A notificação pessoal traz ela de volta ao app toda vez que tem atendimento
Resistência à mudança da agenda de papel
O app imita a lógica da agenda tradicional: simples e sem complexidade
Custo da WhatsApp API afastar usuárias do Pro
O plano gratuito já entrega valor real; a API é um upgrade opcional
Aprovação de templates pela Meta demorar
Lançamento inicial apenas com link manual; API ativada gradualmente


